\section{Where the pieces usually live}

You should not need a forensic map of someone else's laptop to write a chapter. These habits exist so ``where did that paragraph go?'' stays a boring question---and so a colleague who only speaks LilyPond never has to wade through your prose.

\begin{bookStyledList}{Common locations}
\item \textbf{The book's outline:} \verb|main.tex|. Treat it as the driver: document setup, then a list of \verb|\input| lines in reading order. It should orchestrate chapters and back matter, not hold long passages of finished prose.
\item \textbf{Chapter prose:} one section per file at the chapter root is a comfortable default.
\item \textbf{Chapter music only:} keep music in the inline and excerpt areas as standalone \texttt{.ly} files and pull them in with \verb|\lilypondfile{...}| from small \texttt{.tex} chapter fragments so you can run \texttt{make check-ly} for native LilyPond errors. Full-page block examples can stay inline or follow the same pattern when you want the same check.
\item \textbf{Page appearance and overall typography:} shared layout and package setup for teammates who enjoy fine-tuning headings, spacing, and fonts.
\item \textbf{Numbered music examples:} shared LaTeX helpers that introduce each example, add it to the list at the front of the book, and optionally put tall scores in a fixed-height frame so vertical space is shared between systems. A separate pair of helpers handles scores that must break across pages at a chosen system boundary. Section~\ref{sec:book-commands} lists every command name and a minimal nesting example you can paste into a chapter file.
\item \textbf{Chord symbols inside ordinary sentences:} a small chord helper so accidentals and spacing stay consistent when harmony appears in running text.
\item \textbf{Default engraving layout for snippets:} shared LilyPond include files for block examples and for inline excerpts, where staff size and spacing defaults live.
\item \textbf{References and further reading:} the shared bibliography file, which feeds the bibliography at the end of the PDF.
\end{bookStyledList}

Installation, Docker, GitHub Actions, and the exact rebuild incantations live in the repository \texttt{README}. Treat this PDF as the story and the map; treat that file as the checklist you hand to a collaborator who is brave but busy.
